\chapter{Introduction}\label{ch:intro}

\section{Motivation}\label{sec:motivation}

Search algorithms and index structures are the backbone of modern in-memory databases; trie-based
indexes and cache-line-optimized processing are at the centre of this thesis. Since the Adaptive
Radix Tree (ART) by Leis et al.~\cite{leis2013art} in 2013, research has diversified in several
directions: trie-height optimization (HOT~\cite{binna2018hot}), hybrid B\textsuperscript{+}
structures (Masstree~\cite{mao2012masstree}, B\textsuperscript{2}-tree~\cite{schmeisser2022b2tree}),
succinct representations (CoCo-Trie~\cite{boffa2024coco}, SuRF~\cite{zhang2018surf}), and
self-tuning approaches (START~\cite{fent2020start}).

In parallel, a dedicated body of research on cache awareness, prefetching, and allocation strategies
has emerged. Heuristics such as Cache-Sensitive Search Trees, cache-oblivious layouts, and
hierarchical prefetching each address a separately considered aspect of the search-algorithm
architecture. Because code whose behaviour rests on cache-line-optimized execution is so
fundamental, it is sensible to segment the constituent parts of the search-algorithm shell into
axes. This amounts to the backward abstraction of known search-algorithm components into reusable
modules across the most important state-of-the-art papers.

This thesis closes that gap with a framework whose building blocks are treated as interchangeable
design dimensions of a large design space~\cite{idreos2018periodic, idreos2018datacalculator}.
Structurally different search methods --- from tree-like indexes to non-tree-like hash tables, as
well as special properties such as cache lines --- thereby become comparable as points of the same
design space; non-tree-like methods occupy a restricted subset of the dimensions over the same
search-algorithm interface, as if they were the sole fan of a tree root. As a probe for
state-of-the-art membership, PRT-ART (a Probst Redirect Tree / ART hybrid) serves as an abstract,
incomplete assembly of novel search-algorithm components.

\section{Problem Statement}\label{sec:problem}

Existing cache-aware index designs optimize one or a few aspects in isolation and report results
that are hard to compare across publications, because each work uses its own data structures,
workloads, and measurement methodology. The weights of improvements from individual components
remain inseparable: a reported performance gain cannot be unambiguously attributed to a single
component (node type, layout, value placement, prefetching), because the components stay intertwined
in the original designs. There is no common framework that (i)~decomposes such designs into
interchangeable axes, (ii)~reconstructs published algorithms as explicit configurations, and
(iii)~measures both individual axes and whole algorithms on a single, fair basis.

\section{Research Questions}\label{sec:rqs}
From the task description and the supervision meetings arise one main research question frozen with
the supervisor (RQ0) and four measurable sub-questions (RQ1--RQ4); complementary aspects are
attached as \emph{sub-questions}.
\begin{description}
  \item[RQ0 (Main question)] By how much do \emph{active, measurement-driven} cache-aware decisions
  --- page type, memory layout, value placement, and prefetching --- improve the performance of
  trie-based search structures over \emph{passive, statically dimensioned} variants on different CPU
  platforms (hybrid CPUs, Sapphire Rapids)?
  \item[RQ1 (Composition)] How can a search-algorithm design be canonically decomposed into
  orthogonal axes and systematized in a permutation matrix, such that the state-of-the-art methods
  can be reconstructed bias-free as profile configurations --- ideally from their original code ---
  and permuted arbitrarily? \emph{Sub-question:} Which axis couplings (node type $\times$ SIMD
  $\times$ path compression) limit the isolation of individual axes, and how does the same
  axis-constant structure behave with respect to cache-update and cache-line fill behaviour when
  only the ISA/SIMD axis (SSE2/AVX2/AVX-512, NEON/SVE2, RVV) is varied?
  \item[RQ2 (Measurement methodology)] How must a reproducible measurement pipeline with
  hardware-counter collection be designed so that every axis permutation is measured fairly (all
  else held equal, i.e.\ identical compiler and flag basis) across three measurement series (A--C)
  over three granularities (micro-, macro-, and overall benchmarking), and under which
  platform-calibrated criteria does a platform count as successfully verified? \emph{Sub-question:}
  How can contradictory findings in the literature on the optimal node/cache-line size (one cache
  line for CSS/CSB\textsuperscript{+} versus up to sixteen for Hankins/Patel) be re-measured
  bias-free through axis-isolated permutation with all else held equal rather than adopted from
  incompatible original setups?
  \item[RQ3 (Probe comparison)] Does the probe PRT-ART win against the eight Rank-1 SOTA designs
  (ART, HOT, Masstree, CoCo-Trie, START, B\textsuperscript{2}-tree, Wormhole,
  SuRF)\footnote{The \emph{ranks} order the state-of-the-art publications by centrality for this thesis: Rank-1 comprises the ten most important core works --- among them the eight standalone search designs as well as ARTSync (P08) and LOUDS (P09), which enter as organs ---, Rank-2 the cache-aware layout and B\textsuperscript{+} precursors, and Rank-3 the prefetching, synchronization, and TU-Dresden hardware works.}
  and the Rank-2/3 SOTA \emph{not} through asymptotic novelty but through better microarchitectural
  properties (higher cache-line utilization, fewer LLC and dTLB misses, lower branch cost, smaller
  hot search paths) for exact lookup, prefix lookup, and prefix enumeration --- measured under the
  three mandatory measurement series A (probe vs.\ state of the art), B (systematic variation of the
  design decisions), and C (merge/regression old vs.\ new)? \emph{Sub-question:} Which local page
  representation (8-/16-bit point page versus compact page) is cheapest at a given branch point, and
  by which platform-calibrated switching rules --- including the order
  Array$_{256}\to$Array$_{65535}\to$Vector$\langle$u8,u8$\rangle\to$Vector$\langle$u16,u16$\rangle$
  by fill level and search type --- should switching occur adaptively?
  \item[RQ4 (Compile time vs.\ run time)] Which axis properties must be fixed at compile time to
  produce platform-optimized, provenance-verifiable binaries per permutation, and which can or must
  be selected at run time (bit-encoded permutation identifier, adaptive layout selection) without
  violating scientific comparability (identical compiler and flag basis)?
\end{description}

\section{Objectives and Contributions}\label{sec:contributions}

The objective follows the task description in four steps: decompose cache-aware search algorithms
into interchangeable design constituents, reconstruct state-of-the-art designs as explicit
configurations, find/decompose/measure cache-line-aware behaviour both per constituent and for the
whole algorithm on a common CPU basis, and deliver the best overall algorithm --- following an
autonomous heuristic-profile evaluation and heuristic-function modelling --- as a standalone binary.
The code of the thesis acts as the \emph{formal operator} of the cache engine: it loads the probe
PRT-ART, compiles it via metaprogramming against the core constituents of the cache engine, delivers
ABI-stable \emph{living beings} (search-algorithm instances), and drives them through the probe dock
for measurement. Architecturally the model separates three layers --- a genus interface
(SearchAlgorithm, Container, Graph), beneath it living-being subclasses with a fixed axis set,
beneath that the nineteen axes as organs --- and four subsystems (measurement driver,
CacheEngineBuilder, CacheEngine, probe). This is \emph{one} architecture, not a coexistence:
\emph{living being} and \texttt{SearchAlgorithm} name the same entity (metaphor and technical term),
and the anatomy is not a second model but the \emph{body} of this living being --- the fixed
composition of its nineteen axis-organs as a standalone binary with an ABI-stable interface of a
search algorithm --- optionally with or without the measurement observer for debug, measurement, or
release. The root \texttt{ExecutionEngine} carries both living beings and axis-less \emph{viruses}
(pure graph algorithms) as siblings; the layering
CacheEngine\,$\to$\,ExecutionEngine\,$\to$\,SearchEngine is merely the ABI-runtime \emph{view} of
the same living being across the module boundary, not a parallel construct. The
\texttt{ExecutionEngine} supplies, for all genera, the execution and measurement tooling required to
support measurement-driven heuristic optimization --- OS primitives (cache-line-aligned allocation,
NUMA read pinning, coherence-aware writes) and, in experiment mode, the measurement aggregator. The
concept of the axis architecture itself is adopted from prior work by the author
(Comdare\,/\,BEP~Venture~UG, originally for optimizing deduplication procedures in the UltiHash
database); the contribution of this thesis is its transfer to cache-aware search structures and its
measurement-oriented realization. The central contributions are:
\begin{itemize}
  \item A three-layer module interface (CacheEngine $\to$ ExecutionEngine $\to$ SearchEngine) ---
  the ABI-runtime view of the same living being --- with ABI-stable C++23 module interfaces and
  variadically templated hybrid-API specializations (1~parameter $\Rightarrow$ \texttt{std::vector}
  API; 2~parameters $\Rightarrow$ \texttt{std::map} API; $N>2 \Rightarrow$
  \texttt{std::map<K,\,std::tuple<V$_1$\ldots V$_N$>>}).
  \item A nineteen-axis permutation matrix with \texttt{std::variant} building blocks and
  compile-time fallback between the probe stack and the state-of-the-art stack.
  \item Thirty SOTA algorithm profiles as persisted XML configurations (8~Rank-1 + 22~Rank-2/3) that
  the \texttt{ExperimentDriver} automatically transforms into pre-compiled C++23 modules;
  prospectively, the profile set is to cover every analysed publication with at least one algorithm
  profile.
  \item PRT-ART as a probe with eight building-block layers of its own (4+2 pool allocator,
  optimistic lock coupling with reserved blocks, multi-level layout, path-oriented prefetch, density
  tracker, hypothesis metrics H1/H2/H3, inline/external/chain-ref value handles, signaling-bit
  serialization).
  \item A measurement driver with defined-/full-mode distinction, profile-aware workload routing,
  and three granularities --- micro-, macro-, and overall benchmarking (individual constituents,
  interface operations, whole algorithms).
  \item \emph{As objective and outlook:} delivery of the best overall algorithm --- automatically
  selected from the fine-grained measurement results --- as a standalone, ABI-stable binary behind
  the unified search-algorithm interface; the measurement pipeline already provides the required
  ranking, while the automatic selection and binary shipping are envisioned as an extension.
  \item A generic platform probe that determines cache properties (line size, hierarchy depth,
  topology) per target platform and feeds them as fixed quantities into a system-optimized binary;
  the two production machines are measured for this purpose under two operating systems (immutable
  Talos and root Linux with full hardware-counter access).
  \item A bias-free comparison matrix that runs \emph{all} load profiles against \emph{all}
  reconstructed living beings, thereby countering the methodological home-field bias of the
  individual publications.
\end{itemize}

\section{Structure of the Thesis}\label{sec:structure}

Chapter~\ref{ch:fundamentals} introduces the necessary background as definitional classes.
Chapter~\ref{ch:sota} mirrors these in the state of the art as instances: an overview with corpus
clusters, the cache concepts, the axis-wise dissection of the methods, the workload frameworks, the
measurement technologies, and the architectural design --- closing with the research gap.
Chapter~\ref{ch:concept} presents the axis library framework, the layered architecture, the ABI
interface, and the nineteen-axis permutation matrix, including the thematic axis repositories and
their sub-axes. Chapter~\ref{ch:impl} documents the implementation across the three repositories.
Chapter~\ref{ch:methodology} describes the measurement methodology with the three mandatory
measurement series (A--C), the three granularities (micro-/macro-/overall benchmarking), the real
datasets, and the fairness rules. Chapter~\ref{ch:results} presents and discusses the results,
including the hypotheses H1--H3. Chapter~\ref{ch:conclusion} concludes with a summary and an
outlook. The author's own full text is dimensioned for a length of roughly 60 to 80 pages.
